\documentclass{ximera}

\input{../../../preamble.tex}


\definecolor{fvdnavy}{HTML}{071D43}
\definecolor{fvdcyan}{HTML}{20BFC6}
\definecolor{fvdcyanlight}{HTML}{EFFCFB}
\definecolor{fvdmagenta}{HTML}{F10D69}
\definecolor{fvdmagentalight}{HTML}{FFF1F7}
\definecolor{fvdtext}{HTML}{163044}





%=================================================
% Pedagogische omgevingen
% PDF: tcolorbox
% HTML: learning-box
%=================================================

\newenvironment{voorkennis}
{%
    \ifdefined\HCode
        \HCode{%
<section class="learning-box learning-box--prior">
<div class="learning-box__header">
<span class="learning-box__icon" aria-hidden="true"></span>
<span class="learning-box__title">Voorkennis</span>
</div>
<div class="learning-box__content">
        }%
        \begin{itemize}
    \else
       \begin{tcolorbox}[
    enhanced,
    breakable,
    colback=fvdcyanlight,
    colframe=fvdcyan,
    coltext=fvdtext,
    boxrule=0.5pt,
    leftrule=4pt,
    arc=4mm,
    outer arc=4mm,
    left=7mm,
    right=6mm,
    top=5mm,
    bottom=5mm,
    before skip=6mm,
    after skip=6mm,
    title=Voorkennis,
    coltitle=fvdcyan,
    fonttitle=\bfseries\Large,
    attach title to upper={\par\smallskip},
    boxed title style={
        empty,
        left=0mm,
        right=0mm,
        top=0mm,
        bottom=0mm
    },
    overlay={
        \draw[
            fvdcyan,
            line width=0.5pt
        ]
        ([xshift=42mm,yshift=-13mm]frame.north west)
        --
        ([xshift=-7mm,yshift=-13mm]frame.north east);
    }
]
        \begin{itemize}
    \fi
}
{%
    \end{itemize}
    \ifdefined\HCode
        \HCode{%
</div>
</section>
        }%
    \else
        \end{tcolorbox}
    \fi
}


\newenvironment{leerdoelen}
{%
    \ifdefined\HCode
        \HCode{%
<section class="learning-box learning-box--goals">
<div class="learning-box__header">
<span class="learning-box__icon" aria-hidden="true"></span>
<span class="learning-box__title">Leerdoelen</span>
</div>
<div class="learning-box__content">
        }%
        \begin{itemize}
    \else
\begin{tcolorbox}[
    enhanced,
    breakable,
    colback=fvdmagentalight,
    colframe=fvdmagenta,
    coltext=fvdtext,
    boxrule=0.5pt,
    leftrule=4pt,
    arc=4mm,
    outer arc=4mm,
    left=7mm,
    right=6mm,
    top=5mm,
    bottom=5mm,
    before skip=6mm,
    after skip=6mm,
    title=Leerdoelen,
    coltitle=fvdmagenta,
    fonttitle=\bfseries\Large,
    attach title to upper={\par\smallskip},
    boxed title style={
        empty,
        left=0mm,
        right=0mm,
        top=0mm,
        bottom=0mm
    },
    overlay={
        \draw[
            fvdmagenta,
            line width=0.5pt
        ]
        ([xshift=42mm,yshift=-13mm]frame.north west)
        --
        ([xshift=-7mm,yshift=-13mm]frame.north east);
    }
]
        \begin{itemize}
    \fi
}
{%
    \end{itemize}
    \ifdefined\HCode
        \HCode{%
</div>
</section>
        }%
    \else
        \end{tcolorbox}
    \fi
}


\addPrintStyle{../../..}

\title{Oefeningen — De afgeleide functie}
\author{Ingmar Herreman}

\begin{document}

% FVD_CHAPTER_DATA_START
% Automatisch gegenereerd uit index.tex.
% Niet handmatig aanpassen.
\fvdchapterdata{1}{Veranderingen in beeld}{2}
% FVD_CHAPTER_DATA_END














































































\maketitle

\FVDbeginExercises

% ============================================================
% BASIS
% ============================================================



% ------------------------------------------------------------
% OEFENING 1
% ------------------------------------------------------------

\begin{oefening}[Afgeleid getal of afgeleide functie?]{\oefB\ESS}

Gegeven is

\[
f'(x)=4x^3-6x.
\]

\begin{question}
Is $f'$ een getal of een functie?

\begin{multipleChoice}
  \choice{Een getal}
  \choice[correct]{Een functie}
\end{multipleChoice}
\end{question}

\begin{question}
Bereken

\[
f'(0).
\]
\end{question}

\begin{question}
Bereken

\[
f'(2).
\]
\end{question}

\begin{question}
Leg uit wat $f'(2)$ betekent voor de grafiek van $f$.
\end{question}

\begin{oplossing}

$f'$ is een functie: aan elke invoerwaarde $x$ koppelt ze het afgeleid getal
$f'(x)$.

Voor $x=0$:
\[
f'(0)=4\cdot0^3-6\cdot0=0.
\]

Voor $x=2$:
\[
f'(2)=4\cdot2^3-6\cdot2=32-12=20.
\]

Dus
\[
\boxed{f'(2)=20}.
\]

Meetkundig betekent dit dat de raaklijn aan de grafiek van $f$ in het punt
met abscis $2$ richtingscoëfficiënt $20$ heeft. De grafiek stijgt daar dus
sterk.

\end{oplossing}

\end{oefening}


% ------------------------------------------------------------
% OEFENING 2
% ------------------------------------------------------------

\begin{oefening}[Constante functies]{\oefB\ESS}

Bepaal telkens de afgeleide.

\begin{question}

\(
f(x)=7.
\)

\end{question}

\begin{question}

\(
g(x)=-13.
\)

\end{question}

\begin{question}

\(
h(x)=\pi.
\)

\end{question}

\begin{question}
Leg in woorden uit waarom de afgeleide van een constante functie
nul is.
\end{question}

\begin{oplossing}

De afgeleide van een constante functie is nul. Dus:
\[
f'(x)=0,
\qquad
g'(x)=0,
\qquad
h'(x)=0.
\]

Een constante functie heeft een horizontale grafiek. De richtingscoëfficiënt
van elke raaklijn is daarom nul.

\end{oplossing}

\end{oefening}


% ------------------------------------------------------------
% OEFENING 3
% ------------------------------------------------------------

\begin{oefening}[Machtsregel]{\oefB\ESS}

Bepaal telkens de afgeleide.

\begin{question}

\(
f(x)=x^2.
\)

\end{question}

\begin{question}

\(
g(x)=x^5.
\)

\end{question}

\begin{question}

\(
h(x)=x^8.
\)

\end{question}

\begin{question}

\(
p(x)=x.
\)

\end{question}

\begin{oplossing}

Met de machtsregel
\[
(x^n)'=nx^{n-1}
\]
krijgen we:
\[
f'(x)=2x,
\]
\[
g'(x)=5x^4,
\]
\[
h'(x)=8x^7,
\]
en
\[
p'(x)=1.
\]

\end{oplossing}

\end{oefening}


% ------------------------------------------------------------
% OEFENING 4
% ------------------------------------------------------------

\begin{oefening}[Constante factor]{\oefB\ESS}

Bepaal telkens de afgeleide.

\begin{question}

\(
f(x)=4x^3.
\)

\end{question}

\begin{question}

\(
g(x)=-7x^5.
\)

\end{question}

\begin{question}

\(
h(x)=\frac12 x^6.
\)

\end{question}

\begin{question}

\(
p(x)=-\frac34 x^4.
\)

\end{question}

\begin{oplossing}

We gebruiken de machtsregel en behouden de constante factor:
\[
f'(x)=12x^2,
\]
\[
g'(x)=-35x^4,
\]
\[
h'(x)=3x^5,
\]
en
\[
p'(x)=-3x^3.
\]

\end{oplossing}

\end{oefening}


% ------------------------------------------------------------
% OEFENING 5
% ------------------------------------------------------------

\begin{oefening}[Som en verschil]{\oefB\ESS}

Bepaal telkens de afgeleide.

\begin{question}

\(
f(x)=x^4+3x^2.
\)

\end{question}

\begin{question}

\(
g(x)=2x^5-4x^3+x.
\)

\end{question}

\begin{question}

\(
h(x)=-x^6+5x^2-8.
\)

\end{question}

\begin{oplossing}

We leiden elke term afzonderlijk af.

\[
f(x)=x^4+3x^2
\quad\Longrightarrow\quad
\boxed{f'(x)=4x^3+6x}.
\]

\[
g(x)=2x^5-4x^3+x
\quad\Longrightarrow\quad
\boxed{g'(x)=10x^4-12x^2+1}.
\]

\[
h(x)=-x^6+5x^2-8
\quad\Longrightarrow\quad
\boxed{h'(x)=-6x^5+10x}.
\]

\end{oplossing}

\end{oefening}


% ------------------------------------------------------------
% OEFENING 6
% ------------------------------------------------------------

\begin{oefening}[Veeltermfuncties afleiden]{\oefB\ESS}

Bepaal telkens $f'(x)$.

\begin{question}

\(
f(x)=3x^4-5x^2+7x-9.
\)

\end{question}

\begin{question}

\(
f(x)=-2x^5+4x^3-6x+1.
\)

\end{question}

\begin{question}

\(
f(x)=\frac13x^6-\frac12x^4+5x^2.
\)

\end{question}

\begin{question}

\(
f(x)=7x^8-3x^3+2.
\)

\end{question}

\begin{oplossing}

\[
f(x)=3x^4-5x^2+7x-9
\]
geeft
\[
\boxed{f'(x)=12x^3-10x+7}.
\]

\[
f(x)=-2x^5+4x^3-6x+1
\]
geeft
\[
\boxed{f'(x)=-10x^4+12x^2-6}.
\]

\[
f(x)=\frac13x^6-\frac12x^4+5x^2
\]
geeft
\[
\boxed{f'(x)=2x^5-2x^3+10x}.
\]

\[
f(x)=7x^8-3x^3+2
\]
geeft
\[
\boxed{f'(x)=56x^7-9x^2}.
\]

\end{oplossing}

\end{oefening}


% ------------------------------------------------------------
% OEFENING 7
% ------------------------------------------------------------

\begin{oefening}[Eerst vereenvoudigen]{\oefB\ESS}

Werk eerst uit en bepaal daarna de afgeleide.

\begin{question}

\(
f(x)=2x(x^2-3x+4).
\)

\end{question}

\begin{question}

\(
g(x)=(x+1)(x^2-2x+3).
\)

\end{question}

\begin{question}

\(
h(x)=x(x-2)(x+2).
\)

\end{question}

\begin{oplossing}

We werken eerst de producten uit.

\[
\begin{aligned}
f(x)
&=2x(x^2-3x+4)\\
&=2x^3-6x^2+8x,
\end{aligned}
\]
dus
\[
\boxed{f'(x)=6x^2-12x+8}.
\]

\[
\begin{aligned}
g(x)
&=(x+1)(x^2-2x+3)\\
&=x^3-x^2+x+3,
\end{aligned}
\]
dus
\[
\boxed{g'(x)=3x^2-2x+1}.
\]

\[
\begin{aligned}
h(x)
&=x(x-2)(x+2)\\
&=x(x^2-4)\\
&=x^3-4x,
\end{aligned}
\]
dus
\[
\boxed{h'(x)=3x^2-4}.
\]

\end{oplossing}

\end{oefening}


% ------------------------------------------------------------
% OEFENING 8
% ------------------------------------------------------------

\begin{oefening}[Eerste en tweede afgeleide]{\oefB\ESS}

Bepaal telkens $f'(x)$ en $f''(x)$.

\begin{question}

\(
f(x)=x^4-2x^3+5x.
\)

\end{question}

\begin{question}

\(
f(x)=3x^5-4x^2+7.
\)

\end{question}

\begin{question}

\(
f(x)=-2x^6+x^3-5x.
\)

\end{question}

\begin{oplossing}

Voor
\[
f(x)=x^4-2x^3+5x
\]
geldt
\[
\boxed{f'(x)=4x^3-6x^2+5}
\]
en
\[
\boxed{f''(x)=12x^2-12x}.
\]

Voor
\[
f(x)=3x^5-4x^2+7
\]
geldt
\[
\boxed{f'(x)=15x^4-8x}
\]
en
\[
\boxed{f''(x)=60x^3-8}.
\]

Voor
\[
f(x)=-2x^6+x^3-5x
\]
geldt
\[
\boxed{f'(x)=-12x^5+3x^2-5}
\]
en
\[
\boxed{f''(x)=-60x^4+6x}.
\]

\end{oplossing}

\end{oefening}


% ------------------------------------------------------------
% OEFENING 9
% ------------------------------------------------------------

\begin{oefening}[Functiewaarden van de afgeleiden]{\oefB\ESS}

Gegeven:

\[
f(x)=x^3-3x^2+2x+1.
\]

\begin{question}
Bepaal $f'(x)$.
\end{question}

\begin{question}
Bepaal $f''(x)$.
\end{question}

\begin{question}
Bereken

\[
f'(0),\qquad f'(1),\qquad f'(2).
\]

\end{question}

\begin{question}
Bereken

\[
f''(1).
\]

\end{question}

\begin{oplossing}

Gegeven:
\[
f(x)=x^3-3x^2+2x+1.
\]

Dan
\[
\boxed{f'(x)=3x^2-6x+2}
\]
en
\[
\boxed{f''(x)=6x-6}.
\]

Verder:
\[
f'(0)=2,
\]
\[
f'(1)=3-6+2=-1,
\]
\[
f'(2)=12-12+2=2.
\]

Dus
\[
\boxed{f'(0)=2,\qquad f'(1)=-1,\qquad f'(2)=2}.
\]

Ten slotte:
\[
f''(1)=6-6=0.
\]

\end{oplossing}

\end{oefening}


% ------------------------------------------------------------
% OEFENING 10
% ------------------------------------------------------------

\begin{oefening}[Welke afgeleide klopt?]{\oefB\ESS}

Gegeven:

\[
f(x)=2x^5-3x^3+4x-7.
\]

Welke uitdrukking is $f'(x)$?

\begin{multipleChoice}
\choice{$10x^5-9x^3+4$}
\choice{$10x^4-9x^2+4x$}
\choice[correct]{$10x^4-9x^2+4$}
\choice{$2x^4-3x^2+4$}
\end{multipleChoice}

Leg daarna uit waarom minstens twee van de foute antwoorden niet
kunnen kloppen.

\begin{oplossing}

We leiden term per term af:
\[
f'(x)=10x^4-9x^2+4.
\]

Dus het correcte antwoord is
\[
\boxed{10x^4-9x^2+4}.
\]

Bij het eerste foute antwoord is de macht bij het afleiden niet met één
verminderd. Bij het tweede foute antwoord werd de afgeleide van $4x$
verkeerd als $4x$ geschreven, terwijl
\[
(4x)'=4.
\]

\end{oplossing}

\end{oefening}


% ------------------------------------------------------------
% OEFENING 11
% ------------------------------------------------------------

\begin{oefening}[Snelle controle met de graad]{\oefB}

Een veeltermfunctie $f$ heeft graad $7$.

\begin{question}
Welke graad heeft $f'$?
\end{question}

\begin{question}
Welke graad heeft $f''$?
\end{question}

\begin{question}
Welke graad heeft de derde afgeleide $f'''$?
\end{question}

\begin{question}
Na hoeveel keer afleiden krijg je noodzakelijk een constante functie?
\end{question}

\begin{oplossing}

Bij het afleiden van een niet-constante veelterm daalt de graad telkens met
één.

Als $f$ graad $7$ heeft, dan heeft
\[
f'
\]
graad $6$,
\[
f''
\]
graad $5$,
en
\[
f'''
\]
graad $4$.

Na $7$ keer afleiden krijgen we noodzakelijk een constante functie.
Nog één keer afleiden geeft vervolgens de nulfunctie.

\end{oplossing}

\end{oefening}


% ------------------------------------------------------------
% OEFENING 12
% ------------------------------------------------------------

\begin{oefening}[Constanten en verticale verschuivingen]{\oefB\ESS}

Gegeven zijn

\[
f(x)=3x^4-2x^2+5
\]

en

\[
g(x)=3x^4-2x^2-100.
\]

\begin{question}
Bepaal $f'(x)$.
\end{question}

\begin{question}
Bepaal $g'(x)$.
\end{question}

\begin{question}
Vergelijk beide afgeleiden.
\end{question}

\begin{question}
Leg grafisch uit waarom dit resultaat logisch is.
\end{question}

\begin{oplossing}

Voor
\[
f(x)=3x^4-2x^2+5
\]
geldt
\[
f'(x)=12x^3-4x.
\]

Voor
\[
g(x)=3x^4-2x^2-100
\]
geldt eveneens
\[
g'(x)=12x^3-4x.
\]

Dus
\[
\boxed{f'(x)=g'(x)}.
\]

De grafieken van $f$ en $g$ hebben exact dezelfde vorm en verschillen alleen
door een verticale verschuiving. Een verticale verschuiving verandert de
hellingen van de raaklijnen niet.

\end{oplossing}

\end{oefening}



% ------------------------------------------------------------
% OEFENING 13
% ------------------------------------------------------------

\begin{oefening}[Foutenanalyse]{\oefU\ESS}

Een leerling leidt

\[
f(x)=4x^5-3x^2+7
\]

als volgt af:

\[
f'(x)=20x^5-6x+7.
\]

\begin{question}
Zoek alle fouten.
\end{question}

\begin{question}
Geef de correcte afgeleide.
\end{question}

\begin{question}
Noem bij elke fout de afleidingsregel die verkeerd werd toegepast.
\end{question}

\begin{oplossing}

De leerling schrijft
\[
f'(x)=20x^5-6x+7,
\]
maar maakt drie fouten.

Voor
\[
4x^5
\]
moet de macht met één verminderen:
\[
(4x^5)'=20x^4.
\]

Voor
\[
-3x^2
\]
geldt correct:
\[
(-3x^2)'=-6x.
\]

De afgeleide van de constante $7$ is nul:
\[
7'=0.
\]

Daarom:
\[
\boxed{f'(x)=20x^4-6x}.
\]

De verkeerde regels zijn dus: de machtsregel werd bij de eerste term niet
volledig toegepast en de regel voor de afgeleide van een constante werd bij
de laatste term vergeten.

\end{oplossing}

\end{oefening}


% ------------------------------------------------------------
% OEFENING 14
% ------------------------------------------------------------

\begin{oefening}[Is dit mogelijk?]{\oefU\ESS}

Onderzoek telkens of de uitspraak mogelijk is.

\begin{question}
Een veeltermfunctie van graad $4$ heeft een afgeleide van graad $4$.
\end{question}

\begin{question}
Een veeltermfunctie van graad $4$ heeft een tweede afgeleide van
graad $2$.
\end{question}

\begin{question}
Een niet-constante veeltermfunctie heeft afgeleide nul voor alle $x$.
\end{question}

\begin{question}
Twee verschillende veeltermfuncties kunnen dezelfde afgeleide hebben.
\end{question}

Motiveer elk antwoord.

\begin{oplossing}

\textbf{1. Een veelterm van graad $4$ met afgeleide van graad $4$:}
niet mogelijk.

Bij het afleiden daalt de graad van een niet-constante veelterm met één.
De afgeleide heeft dus graad $3$.

\medskip

\textbf{2. Een veelterm van graad $4$ met tweede afgeleide van graad $2$:}
wel mogelijk.

Na twee keer afleiden daalt de graad van $4$ naar $2$.

\medskip

\textbf{3. Een niet-constante veelterm met afgeleide nul voor alle $x$:}
niet mogelijk.

Alleen een constante veelterm heeft overal afgeleide nul.

\medskip

\textbf{4. Twee verschillende veeltermfuncties met dezelfde afgeleide:}
wel mogelijk.

Bijvoorbeeld
\[
f(x)=x^2
\qquad\text{en}\qquad
g(x)=x^2+5.
\]
Dan
\[
f'(x)=g'(x)=2x.
\]

\end{oplossing}

\end{oefening}


% ------------------------------------------------------------
% OEFENING 15
% ------------------------------------------------------------

\begin{oefening}[Zelfde afgeleide]{\oefU\ESS}

Gegeven:

\[
f(x)=2x^3-5x.
\]

Zoek drie verschillende functies $g$ waarvoor

\[
g'(x)=f'(x).
\]

\begin{question}
Geef drie mogelijke voorschriften.
\end{question}

\begin{question}
Wat hebben deze functies gemeenschappelijk?
\end{question}

\begin{question}
Waarin verschillen hun grafieken?
\end{question}

\begin{question}
Formuleer een algemeen besluit.
\end{question}

\begin{oplossing}

Uit
\[
f(x)=2x^3-5x
\]
volgt
\[
f'(x)=6x^2-5.
\]

Drie mogelijke functies met dezelfde afgeleide zijn bijvoorbeeld
\[
g_1(x)=2x^3-5x+1,
\]
\[
g_2(x)=2x^3-5x-4,
\]
en
\[
g_3(x)=2x^3-5x+10.
\]

Ze hebben allemaal dezelfde afgeleide
\[
6x^2-5.
\]

Hun grafieken hebben dus dezelfde vorm en op elke abscis dezelfde helling,
maar ze zijn verticaal ten opzichte van elkaar verschoven.

Algemeen geldt: functies die alleen verschillen door een constante hebben
dezelfde afgeleide.

\end{oplossing}

\end{oefening}


% ------------------------------------------------------------
% OEFENING 16
% ------------------------------------------------------------

\begin{oefening}[Van $f'$ terug naar $f$]{\oefU\ESS}

Van een veeltermfunctie $f$ is gegeven:

\[
f'(x)=6x^2-4x+3.
\]

\begin{question}
Geef één mogelijke functie $f$.
\end{question}

\begin{question}
Geef een tweede mogelijke functie $f$.
\end{question}

\begin{question}
Waarom kun je $f$ niet uniek bepalen met alleen deze informatie?
\end{question}

\begin{question}
Welke extra informatie zou voldoende zijn om één specifieke functie
te bepalen?
\end{question}

\begin{oplossing}

We zoeken een functie waarvan
\[
f'(x)=6x^2-4x+3.
\]

Een mogelijke functie is
\[
\boxed{f(x)=2x^3-2x^2+3x}.
\]

Een tweede mogelijkheid is bijvoorbeeld
\[
\boxed{f(x)=2x^3-2x^2+3x+5}.
\]

De functie is niet uniek bepaald omdat een willekeurige constante bij het
afleiden verdwijnt.

Algemeen:
\[
f(x)=2x^3-2x^2+3x+C.
\]

Eén extra functiewaarde, bijvoorbeeld
\[
f(0)=7,
\]
zou voldoende zijn om de constante $C$ en dus één specifieke functie te
bepalen.

\end{oplossing}

\end{oefening}


% ------------------------------------------------------------
% OEFENING 17
% ------------------------------------------------------------

\begin{oefening}[Grafieken koppelen]{\oefU\ESS}

De drie grafieken hieronder stellen een functie, haar eerste afgeleide
en haar tweede afgeleide voor.

\begin{center}
\begin{tikzpicture}[x=0.85cm,y=0.7cm]

  % grafiek A
  \begin{scope}[xshift=-5.2cm]
    \draw[->,draw=black!45] (-2.5,0)--(2.5,0);
    \draw[->,draw=black!45] (0,-2.4)--(0,3.3);
    \draw[draw=fvdnavy,very thick,domain=-2:2,samples=100]
      plot (\x,{0.4*(\x)^3-0.8*\x});
    \node at (0,-2.8) {A};
  \end{scope}

  % grafiek B
  \begin{scope}
    \draw[->,draw=black!45] (-2.5,0)--(2.5,0);
    \draw[->,draw=black!45] (0,-2.4)--(0,3.3);
    \draw[draw=fvdnavy,very thick,domain=-2:2,samples=100]
      plot (\x,{1.2*(\x)^2-0.8});
    \node at (0,-2.8) {B};
  \end{scope}

  % grafiek C
  \begin{scope}[xshift=5.2cm]
    \draw[->,draw=black!45] (-2.5,0)--(2.5,0);
    \draw[->,draw=black!45] (0,-2.4)--(0,3.3);
    \draw[draw=fvdnavy,very thick,domain=-2:2,samples=100]
      plot (\x,{2.4*\x});
    \node at (0,-2.8) {C};
  \end{scope}

\end{tikzpicture}
\end{center}

\begin{question}
Welke grafiek kan $f$ voorstellen?
\end{question}

\begin{question}
Welke grafiek kan $f'$ voorstellen?
\end{question}

\begin{question}
Welke grafiek kan $f''$ voorstellen?
\end{question}

\begin{question}
Motiveer je keuze met behulp van de graden en de vorm van de grafieken.
\end{question}

\begin{oplossing}

Grafiek A is kubisch, grafiek B is kwadratisch en grafiek C is lineair.

Omdat bij het afleiden van een veelterm de graad telkens met één daalt,
past de volgorde:
\[
\boxed{f:\ \text{A}},
\qquad
\boxed{f':\ \text{B}},
\qquad
\boxed{f'':\ \text{C}}.
\]

Dit klopt ook inhoudelijk: de nulpunten van grafiek B komen overeen met
horizontale raaklijnen van grafiek A, en de lineaire grafiek C beschrijft
hoe de kwadratische grafiek B verandert.

\end{oplossing}

\end{oefening}


% ------------------------------------------------------------
% OEFENING 18
% ------------------------------------------------------------

\begin{oefening}[Van grafiek naar teken van $f'$]{\oefU\ESS}

Een functie $f$ heeft onderstaande grafiek.

\begin{center}
\begin{tikzpicture}[x=1cm,y=0.75cm]

  \draw[->,draw=black!45] (-4,0)--(4.3,0)
    node[right] {$x$};
  \draw[->,draw=black!45] (0,-3)--(0,4.2)
    node[above] {$y$};

  \draw[
    draw=fvdnavy,
    very thick,
    smooth,
    domain=-3.2:3.2,
    samples=150
  ]
    plot (\x,{0.16*(\x)^3-1.2*\x});

  \draw[dashed,draw=black!40] (-1.58,0)--(-1.58,1.26);
  \draw[dashed,draw=black!40] (1.58,0)--(1.58,-1.26);

\end{tikzpicture}
\end{center}

\begin{question}
Op welke delen van de grafiek verwacht je $f'(x)>0$?
\end{question}

\begin{question}
Waar verwacht je $f'(x)<0$?
\end{question}

\begin{question}
Waar verwacht je $f'(x)=0$?
\end{question}

\begin{question}
Schets kwalitatief een mogelijke grafiek van $f'$.
\end{question}

\begin{oplossing}

De grafiek heeft horizontale raaklijnen rond
\[
x\approx -1.58
\qquad\text{en}\qquad
x\approx 1.58.
\]

Links van $-1.58$ stijgt $f$, dus
\[
f'(x)>0.
\]

Tussen $-1.58$ en $1.58$ daalt $f$, dus
\[
f'(x)<0.
\]

Rechts van $1.58$ stijgt $f$ opnieuw, dus
\[
f'(x)>0.
\]

Bij
\[
x\approx -1.58
\qquad\text{en}\qquad
x\approx 1.58
\]
geldt
\[
f'(x)=0.
\]

Een kwalitatieve grafiek van $f'$ is dus een naar boven geopende parabool
met nulpunten rond $-1.58$ en $1.58$.

\end{oplossing}

\end{oefening}


% ------------------------------------------------------------
% OEFENING 19
% ------------------------------------------------------------

\begin{oefening}[Eerste en tweede afgeleide interpreteren]{\oefU\ESS}

Gegeven is een functie $f$ waarvoor op een bepaald punt $x=a$ geldt:

\[
f'(a)=4
\]

en

\[
f''(a)=-3.
\]

\begin{question}
Wat zegt $f'(a)=4$ over de raaklijn aan de grafiek van $f$?
\end{question}

\begin{question}
Wat zegt het teken van $f''(a)$ over de verandering van $f'$ in de
buurt van $a$?
\end{question}

\begin{question}
Kun je uit deze twee waarden alleen besluiten dat $f$ in $a$ een
maximum of minimum heeft? Motiveer.
\end{question}

\begin{oplossing}

Uit
\[
f'(a)=4
\]
volgt dat de raaklijn aan de grafiek van $f$ in $x=a$
richtingscoëfficiënt $4$ heeft. De grafiek stijgt daar.

Verder geldt
\[
f''(a)=-3.
\]
Dit betekent dat $f'$ op dat ogenblik afneemt: de grafiek van $f'$ heeft
bij $x=a$ een negatieve helling.

We kunnen hieruit geen maximum of minimum van $f$ in $a$ besluiten.
Integendeel: omdat
\[
f'(a)=4\neq0,
\]
kan $a$ voor een afleidbare functie geen lokaal extremum zijn.

\end{oplossing}

\end{oefening}


% ------------------------------------------------------------
% OEFENING 20
% ------------------------------------------------------------

\begin{oefening}[Positie, snelheid en versnelling]{\oefU\ESS}

De positie van een bewegend voorwerp wordt gegeven door

\[
s(t)=t^3-6t^2+9t+2,
\]

met $s$ in meter en $t$ in seconde.

\begin{question}
Bepaal de snelheidsfunctie

\[
v(t)=s'(t).
\]
\end{question}

\begin{question}
Bepaal de versnellingsfunctie

\[
a(t)=s''(t).
\]
\end{question}

\begin{question}
Bereken $v(1)$.
\end{question}

\begin{question}
Bereken $a(1)$.
\end{question}

\begin{question}
Geef bij beide resultaten de juiste eenheid.
\end{question}

\begin{question}
Leg uit waarom de eenheden van $s$, $v$ en $a$ verschillend zijn.
\end{question}

\begin{oplossing}

De snelheid is de eerste afgeleide van de positie:
\[
v(t)=s'(t)=3t^2-12t+9.
\]

De versnelling is de afgeleide van de snelheid:
\[
a(t)=s''(t)=6t-12.
\]

Voor $t=1$:
\[
v(1)=3-12+9=0.
\]

Dus
\[
\boxed{v(1)=0\ \mathrm{m/s}}.
\]

Verder:
\[
a(1)=6-12=-6.
\]

Dus
\[
\boxed{a(1)=-6\ \mathrm{m/s^2}}.
\]

De positie heeft als eenheid meter, de snelheid meter per seconde en de
versnelling meter per seconde per seconde. Elke afleiding naar de tijd voegt
dus een factor ``per seconde'' toe.

\end{oplossing}

\end{oefening}


% ------------------------------------------------------------
% OEFENING 21
% ------------------------------------------------------------

\begin{oefening}[Wetenschappelijke eenheden]{\oefU}

De concentratie van een stof wordt beschreven door $C(t)$,
uitgedrukt in $\mathrm{mol/L}$, met $t$ in seconde.

\begin{question}
Welke eenheid heeft $C'(t)$?
\end{question}

\begin{question}
Welke eenheid heeft $C''(t)$?
\end{question}

\begin{question}
Geef een mogelijke betekenis aan een negatieve waarde van $C'(t)$.
\end{question}

\begin{question}
Wat betekent het wanneer $C'(t)<0$ maar $C''(t)>0$?
Beschrijf dit kwalitatief.
\end{question}

\begin{oplossing}

Omdat $C$ wordt uitgedrukt in $\mathrm{mol/L}$ en $t$ in seconde, heeft
\[
C'(t)
\]
als eenheid
\[
\boxed{\mathrm{mol/(L\,s)}}.
\]

Voor de tweede afgeleide:
\[
\boxed{\mathrm{mol/(L\,s^2)}}.
\]

Als
\[
C'(t)<0,
\]
dan neemt de concentratie op dat ogenblik af.

Als tegelijk
\[
C'(t)<0
\qquad\text{en}\qquad
C''(t)>0,
\]
dan daalt de concentratie nog steeds, maar wordt $C'(t)$ groter.
De daling wordt dus minder sterk: de concentratie neemt af, maar steeds
trager.

\end{oplossing}

\end{oefening}


% ============================================================
% GEVORDERD
% ============================================================



% ------------------------------------------------------------
% OEFENING 22
% ------------------------------------------------------------

\begin{oefening}[Een onbekende parameter]{\oefG}

Gegeven is

\[
f(x)=ax^4+bx^2+3x-1.
\]

Er geldt:

\[
f'(1)=9
\]

en

\[
f''(1)=18.
\]

Bepaal $a$ en $b$.

\begin{oplossing}

Gegeven:
\[
f(x)=ax^4+bx^2+3x-1.
\]

Dan
\[
f'(x)=4ax^3+2bx+3
\]
en
\[
f''(x)=12ax^2+2b.
\]

Uit
\[
f'(1)=9
\]
volgt
\[
4a+2b+3=9,
\]
dus
\[
2a+b=3.
\]

Uit
\[
f''(1)=18
\]
volgt
\[
12a+2b=18,
\]
dus
\[
6a+b=9.
\]

Aftrekken geeft
\[
4a=6,
\]
zodat
\[
a=\frac32.
\]

Daaruit:
\[
2\cdot\frac32+b=3
\]
en dus
\[
b=0.
\]

Bijgevolg
\[
\boxed{a=\frac32,\qquad b=0}.
\]

\end{oplossing}

\end{oefening}


% ------------------------------------------------------------
% OEFENING 23
% ------------------------------------------------------------

\begin{oefening}[Welke functie kan dit zijn?]{\oefG}

Een veeltermfunctie $f$ voldoet aan

\[
f''(x)=12x-6.
\]

Bovendien:

\[
f'(0)=4
\]

en

\[
f(0)=2.
\]

Bepaal $f(x)$.

\begin{oplossing}

Gegeven:
\[
f''(x)=12x-6.
\]

Een functie met deze afgeleide is
\[
f'(x)=6x^2-6x+C.
\]

Omdat
\[
f'(0)=4,
\]
volgt
\[
C=4.
\]

Dus
\[
f'(x)=6x^2-6x+4.
\]

Daaruit volgt
\[
f(x)=2x^3-3x^2+4x+D.
\]

Omdat
\[
f(0)=2,
\]
is
\[
D=2.
\]

Bijgevolg:
\[
\boxed{f(x)=2x^3-3x^2+4x+2}.
\]

\end{oplossing}

\end{oefening}


% ------------------------------------------------------------
% OEFENING 24
% ------------------------------------------------------------

\begin{oefening}[Een robotarm]{\oefG}

De positie van het uiteinde van een robotarm langs één as wordt
gemodelleerd door

\[
s(t)=\frac12t^4-4t^3+9t^2+2.
\]

\begin{question}
Bepaal $v(t)$ en $a(t)$.
\end{question}

\begin{question}
Bepaal de snelheid en versnelling op $t=1$.
\end{question}

\begin{question}
Op welke tijdstippen is de snelheid nul?
\end{question}

\begin{question}
Onderzoek op die tijdstippen ook de versnelling.
\end{question}

\begin{question}
Welke bijkomende informatie zou je nodig hebben om te beslissen of
de robotarm daar van bewegingsrichting verandert?
\end{question}

\begin{oplossing}

We leiden de positie af:
\[
v(t)=s'(t)=2t^3-12t^2+18t.
\]

We kunnen factoriseren:
\[
\boxed{v(t)=2t(t-3)^2}.
\]

De versnelling is
\[
a(t)=v'(t)=6t^2-24t+18.
\]

Voor $t=1$:
\[
v(1)=2-12+18=8,
\]
dus
\[
\boxed{v(1)=8}.
\]

Verder:
\[
a(1)=6-24+18=0,
\]
dus
\[
\boxed{a(1)=0}.
\]

De snelheid is nul als
\[
2t(t-3)^2=0.
\]
Daaruit:
\[
\boxed{t=0\quad\text{of}\quad t=3}.
\]

Op deze tijdstippen:
\[
a(0)=18
\]
en
\[
a(3)=54-72+18=0.
\]

Om te beslissen of de robotarm op zo'n tijdstip van bewegingsrichting
verandert, moeten we het teken van $v(t)$ vlak vóór en vlak na dat tijdstip
onderzoeken. Alleen weten dat $v(t)=0$ is daarvoor niet voldoende.

\end{oplossing}

\end{oefening}


% ------------------------------------------------------------
% OEFENING 25
% ------------------------------------------------------------

\begin{oefening}[Een redenering beoordelen]{\oefG}

Een leerling beweert:

\begin{quote}
Als $f''(x)=0$, dan moet $f'(x)$ constant zijn.
\end{quote}

\begin{question}
Is deze uitspraak correct zoals ze geformuleerd is?
\end{question}

\begin{question}
Leg het verschil uit tussen

\[
f''(a)=0
\]

voor één waarde $a$, en

\[
f''(x)=0
\]

voor alle $x$ in een interval.
\end{question}

\begin{question}
Geef voor beide situaties een voorbeeld.
\end{question}

\begin{oplossing}

De uitspraak is niet correct zoals ze geformuleerd is.

Als
\[
f''(a)=0
\]
voor één bepaalde waarde $a$, dan volgt daar niet uit dat $f'$ constant is.

Bijvoorbeeld:
\[
f(x)=x^3.
\]
Dan
\[
f'(x)=3x^2
\]
en
\[
f''(x)=6x.
\]
Dus
\[
f''(0)=0,
\]
maar $f'$ is duidelijk niet constant.

Als daarentegen
\[
f''(x)=0
\]
voor alle $x$ in een interval, dan is $f'$ op dat interval wel constant.

Bijvoorbeeld:
\[
f(x)=3x+2.
\]
Dan
\[
f'(x)=3
\]
en
\[
f''(x)=0
\]
voor alle $x$.

\end{oplossing}

\end{oefening}


% ------------------------------------------------------------
% OEFENING 26
% ------------------------------------------------------------

\begin{oefening}[Construeer zelf een functie]{\oefG}

Construeer een veeltermfunctie $f$ die voldoet aan:

\[
f'(0)=2,
\]

\[
f''(0)=-4,
\]

en

\[
f(0)=3.
\]

\begin{question}
Geef één mogelijke functie.
\end{question}

\begin{question}
Geef een tweede, verschillende functie die aan dezelfde voorwaarden
voldoet.
\end{question}

\begin{question}
Leg uit waarom de voorwaarden de functie niet uniek bepalen.
\end{question}

\begin{oplossing}

We zoeken een veelterm waarvoor
\[
f(0)=3,\qquad f'(0)=2,\qquad f''(0)=-4.
\]

Een eenvoudige keuze is
\[
\boxed{f(x)=3+2x-2x^2}.
\]

Inderdaad:
\[
f'(x)=2-4x
\]
en
\[
f''(x)=-4.
\]
Dus
\[
f'(0)=2,
\qquad
f''(0)=-4,
\qquad
f(0)=3.
\]

Een tweede mogelijkheid is bijvoorbeeld
\[
\boxed{g(x)=3+2x-2x^2+x^3}.
\]

Dan
\[
g'(x)=2-4x+3x^2
\]
en
\[
g''(x)=-4+6x,
\]
zodat opnieuw
\[
g'(0)=2,\qquad g''(0)=-4,\qquad g(0)=3.
\]

De voorwaarden leggen alleen bepaalde waarden van de functie en haar
afgeleiden in $x=0$ vast. Hogeregraadstermen kunnen worden toegevoegd
zonder deze drie voorwaarden noodzakelijk te veranderen.

\end{oplossing}

\end{oefening}





\begin{oefening}[Alles samen]{\oefU\ESS}

Gegeven is

\[
f(x)=2x^5-5x^3+3x^2-7.
\]

\begin{question}
Bepaal $f'(x)$.
\end{question}

\begin{question}
Bepaal $f''(x)$.
\end{question}

\begin{question}
Bereken

\[
f'(1).
\]
\end{question}

\begin{question}
Interpreteer $f'(1)$ meetkundig.
\end{question}

\begin{question}
Bereken

\[
f''(1).
\]
\end{question}

\begin{question}
Leg in woorden uit welk verschil er is tussen de informatie in
$f'(1)$ en de informatie in $f''(1)$.
\end{question}

\begin{question}
Een leerling schrijft:

\[
f'(x)=10x^4-15x^2+6x-7.
\]

Leg zonder alles opnieuw te berekenen uit waarom deze afgeleide
zeker fout is.
\end{question}

\begin{oplossing}

Gegeven:
\[
f(x)=2x^5-5x^3+3x^2-7.
\]

De eerste afgeleide is
\[
\boxed{f'(x)=10x^4-15x^2+6x}.
\]

De tweede afgeleide is
\[
\boxed{f''(x)=40x^3-30x+6}.
\]

Voor $x=1$:
\[
f'(1)=10-15+6=1.
\]

Dus
\[
\boxed{f'(1)=1}.
\]

Meetkundig betekent dit dat de raaklijn aan de grafiek van $f$ in
$x=1$ richtingscoëfficiënt $1$ heeft.

Verder:
\[
f''(1)=40-30+6=16.
\]

Dus
\[
\boxed{f''(1)=16}.
\]

$f'(1)$ beschrijft de helling van de grafiek van $f$ in $x=1$.
$f''(1)$ beschrijft hoe die helling daar verandert.

De voorgestelde afgeleide
\[
10x^4-15x^2+6x-7
\]
kan zeker niet kloppen, omdat de constante term $-7$ bij het afleiden moet
verdwijnen:
\[
(-7)'=0.
\]

\end{oplossing}

\end{oefening}



\FVDendExercises
\end{document}